% Example 4 (fragmentation, clustered): the easy steps form a contiguous run and are
% chained -- automated up to the endpoint, which is augmented -- and the box groups them
% into a single AI-chain task. The hard steps that follow are performed manually. Colors
% follow Figures 1-2 (manual=gray, automated=green, augmented=orange; AI chain = blue box).
\begin{center}
\begin{adjustbox}{max width=\linewidth}
\begin{tikzpicture}[>=latex,
  cbox/.style={rectangle, rounded corners, draw, minimum width=2.0cm, minimum height=0.95cm, align=center, font=\small},
  manual/.style={fill=gray!10},
  automated/.style={fill=green!30},
  augmented/.style={fill=orange!30},
  sub/.style={draw=none, font=\scriptsize, text=black!70},
  chainbox/.style={draw=blue, rectangle, rounded corners, dashed, inner sep=0.22cm, line width=1pt}]

  \node[cbox, automated] (s1) {Step 1\\(Automated)};
  \node[cbox, automated, right=0.7cm of s1] (s2) {Step 2\\(Automated)};
  \node[draw=none, right=0.45cm of s2] (dots1) {$\cdots$};
  \node[cbox, augmented, right=0.45cm of dots1] (sk) {Step $\tfrac{m}{2}$\\(Augmented)};
  \node[cbox, manual, right=0.95cm of sk] (sk1) {Step $\tfrac{m}{2}{+}1$\\(Manual)};
  \node[draw=none, right=0.45cm of sk1] (dots2) {$\cdots$};
  \node[cbox, manual, right=0.45cm of dots2] (sm) {Step $m$\\(Manual)};

  \node[sub, below=1pt of s1] {easy};
  \node[sub, below=1pt of s2] {easy};
  \node[sub, below=1pt of sk] {easy};
  \node[sub, below=1pt of sk1] {hard};
  \node[sub, below=1pt of sm] {hard};

  \draw[->,thick] (s1)--(s2);
  \draw[->,thick] (s2)--(dots1);
  \draw[->,thick] (dots1)--(sk);
  \draw[->,thick] (sk)--(sk1);
  \draw[->,thick] (sk1)--(dots2);
  \draw[->,thick] (dots2)--(sm);

  \node[chainbox, fit=(s1)(s2)(dots1)(sk)] (chain) {};
  \node[font=\footnotesize, text=blue!60!black] at ($(chain.south)+(0,-0.34)$) {single task (AI chain)};
\end{tikzpicture}
\end{adjustbox}
\end{center}
